
\section{ТЕОРЕТИЧЕСКОЕ ОБОСНОВАНИЕ РЕШЕНИЯ}

В главе формально описывается разрабатываемая система и обосновываются принятые архитектурные решения. Сначала приводится формальная постановка задачи и логика работы программного обеспечения (2.1). Затем рассматриваются архитектура и стек технологий, обоснование выбора каждой компоненты (2.2). В конце разбирается прохождение типового запроса через систему (2.3).

\subsection{Логика работы разрабатываемого программного обеспечения}

\subsubsection{Формальная постановка задачи обогащения}

Пусть товарная позиция, поступающая в систему от партнёра, представлена кортежем партнёр-доступных полей

\begin{equation}
\label{eq:2.1}
x = (x^{\text{name}}, x^{\text{brand}}, x^{\text{ingr}}, x^{\text{qty}}),
\end{equation}

где компоненты — название, бренд, состав и количество соответственно; пустым может быть любое поле, кроме названия. Каждому товару ставится в соответствие категория $c \in \mathcal{C} = \{c_1, \dots, c_M\} \cup \{\text{OOD}\}$ ($M = 3$: pasta, chocolate, cheeses; OOD — выход за пределы поддерживаемых доменов), определяемая отдельной моделью

\begin{equation}
\label{eq:2.2}
c = f_0(x),
\end{equation}

(Слой 0, см. 2.2.1). Для каждой категории задана схема $\mathcal{A}_c = \{a_1, \dots, a_{K_c}\}$ с областями $\mathcal{Y}_k$; система строит

\begin{equation}
\label{eq:2.3}
\hat{y} = (\hat{y}_1, \dots, \hat{y}_{K_c}) \in (\mathcal{Y}_1 \cup \{\text{null}\}) \times \dots \times (\mathcal{Y}_{K_c} \cup \{\text{null}\}),
\end{equation}

где $\hat{y}_k = \text{null}$ — явный отказ от ответа. Честное молчание лучше неверного заполнения: повторное обращение к оператору дешевле ошибки на витрине.

\subsubsection{Каскадная схема принятия решения}

Логика обогащения реализована как каскад из четырёх содержательных слоёв; ему предшествует определение категории. Для каждого атрибута $k$ итоговое предсказание формируется по правилу

\begin{gather}
\label{eq:2.4}
\hat{y}_k =
\begin{cases}
f_k^{\text{regex}}(x), & f_k^{\text{regex}}(x) \neq \varnothing; \\[4pt]
f_k^{\text{ML}}(x), & f_k^{\text{regex}}(x) = \varnothing,\ P_{\text{ML}}(\hat{y}_k \mid x) \geq \tau_k,\ P_B(\hat{y}_k \mid \mathbf{e}) \geq \theta_k; \\[4pt]
f_k^{\text{LLM}}(x), & \text{otherwise}.
\end{cases}
\end{gather}

Здесь $f_k^{\text{regex}}$ — детерминированный экстрактор регулярных выражений (Слой 1); $f_k^{\text{ML}}$ — классификатор машинного обучения (Слой 2) над многоязычным векторным представлением партнёр-доступных полей с калиброванной апостериорной вероятностью $P_{\text{ML}}$ и поатрибутным порогом $\tau_k$; $P_B(\hat{y}_k \mid \mathbf{e})$ — оценка байесовской сети при свидетельствах $\mathbf{e} = (x^{\text{brand}}, \hat{y}_{j \neq k})$ с порогом $\theta_k$, заданным как пятый перцентиль распределения $P_B$ по верно размеченным обучающим примерам; $f_k^{\text{LLM}}$ — запасной слой на большой языковой модели (LLM, large language model), возвращающий значение из $\mathcal{Y}_k$ или $\text{null}$. Предсказание Слоя 2 принимается тогда и только тогда, когда одновременно выполнены

\begin{equation}
\label{eq:2.5}
\bigl[P_{\text{ML}}(\hat{y}_k \mid x) \geq \tau_k\bigr] \ \wedge\ \bigl[P_B(\hat{y}_k \mid \mathbf{e}) \geq \theta_k\bigr].
\end{equation}

Это конъюнктивное условие формализует требование темы об ансамблевой оценке уверенности: не параллельное голосование, а последовательная фильтрация с правом каждого слоя отказаться. По итогам 3.3.4 валидатор включён точечно — только на парах «категория — атрибут» с неотрицательным приростом. Схема даёт ступенчатое распределение нагрузки (дешёвый Слой 1 — массовые случаи, дорогой Слой 4 — трудный остаток), отбраковку статистически маловероятных пар «бренд × значение» и архитектурную объяснимость: для каждой ячейки известно, какой слой принял решение.

\subsubsection{Поведение системы при выходе категории за пределы поддерживаемых}

Слой 0 определяет $c = f_0(x)$ до запуска каскада. Если уверенность ниже порога OOD-детектора, товару присваивается метка OOD и система возвращает «категория не поддерживается»: $f_0(x) \in \mathcal{C} \setminus \{\text{OOD}\} \Rightarrow$ запуск per-category каскада со схемой $\mathcal{A}_{f_0(x)}$; $f_0(x) = \text{OOD} \Rightarrow \hat{y} = \varnothing$. Явный отказ выбран по соображениям эксплуатационной безопасности: если товар ошибочно отнести к чужой категории, всем атрибутам достанется неверная схема.

\subsubsection{Метрики качества и стоимости}

Каскад поддерживает явный отказ от ответа; используются две согласованные метрики качества и одна сводная по стоимости.

\begin{equation}
\label{eq:2.6}
\text{acc}_k^{\text{cov}} = \frac{\bigl|\{i : \hat{y}_k^{(i)} = y_k^{(i)} \ \wedge\ \hat{y}_k^{(i)} \neq \text{null}\}\bigr|}{\bigl|\{i : \hat{y}_k^{(i)} \neq \text{null}\}\bigr|}.
\end{equation}

Точность на покрытой части $\text{acc}_k^{\text{cov}}$ — доля верных предсказаний среди ячеек с непустым ответом.

\begin{equation}
\label{eq:2.7}
\text{cov}_k = \frac{\bigl|\{i : \hat{y}_k^{(i)} \neq \text{null}\}\bigr|}{N},
\end{equation}

Покрытие $\text{cov}_k$ — доля ячеек, на которых система решилась дать ответ (где $N$ — мощность проверочной выборки).

\begin{equation}
\label{eq:2.8}
\text{acc}_k^{\text{e2e}} = \text{acc}_k^{\text{cov}} \cdot \text{cov}_k.
\end{equation}

Сквозная точность $\text{acc}_k^{\text{e2e}}$ — основная метрика и заголовочное число (3.3.2); ячейки с $\hat{y}_k^{(i)} = \text{null}$ засчитываются как неверные, что не даёт системе завысить качество за счёт молчания.

\begin{equation}
\label{eq:2.9}
\text{cost}_{\text{LLM}} = \frac{1000}{N \cdot K_c} \sum_{i=1}^{N} \sum_{k=1}^{K_c} \mathbb{1}\bigl[\hat{y}_k^{(i)} \text{ получен через } f_k^{\text{LLM}}\bigr].
\end{equation}

Стоимость $\text{cost}_{\text{LLM}}$ — ожидаемая доля вызовов запасного слоя на 1000 ячеек; метрика инвариантна к выбору модели Слоя 4 и характеризует структуру каскада. Совместный анализ $(\text{acc}^{\text{e2e}}, \text{cost}_{\text{LLM}})$ даёт матрицу «точность — стоимость», на которой строится заголовочный результат (3.3.2).

\subsection{Архитектура и стек технологий}

В разделе описывается архитектурное решение, реализующее формальную постановку 2.1, и обосновывается выбор каждой компоненты стека.

\subsubsection{Архитектура системы}

Функциональная модель системы в виде диаграммы потоков данных приведена на рисунке 2.1. Внешние акторы (партнёр-продавец и витрина каталога интернет-магазина) обмениваются данными со шлюзом API и ML-сервисом, который содержит каскад из пяти слоёв и обращается к хранилищам моделей XGBoost, байесовских сетей и внешнему API большой языковой модели.

\begin{figure}
  \includegraphics[width=\textwidth]{fig_2_1_functional_model.png}
  \caption{Функциональная модель системы обогащения товарных данных}
  \label{fig:fig_2_1_functional_model}
\end{figure}

Система реализована как пятислойный конвейер (Слой 0 + четыре содержательных слоя). Состав и роль каждого слоя:

\begin{itemize}
    \item \textbf{Слой 0 — категорайзер.} LightGBM поверх TF-IDF-представления (term frequency × inverse document frequency, частотно-инверсный показатель) названия и бренда с порогом OOD; определяет $c \in \mathcal{C}$ либо относит к OOD.
    \item \textbf{Слой 1 — экстрактор регулярных выражений.} Извлекает значения, явно присутствующие в тексте (форма пасты, \% какао, цельнозерновой, шаблоны количества). Применяется только к шаблонам с высокой точностью совпадения.
    \item \textbf{Слой 2 — классификатор машинного обучения.} Основной слой: партнёр-доступные поля конкатенируются и кодируются SBERT в 768-мерный вектор (модель paraphrase-multilingual-mpnet-base-v2, архитектура MPNet); поатрибутный XGBoost с порогом $\tau_k$ на отложенной выборке. Изотоническая пост-калибровка исследована в двух режимах (3.3.3.3): серебряная деградирует на эталоне из-за дрейфа; перекрёстная на эталонной выборке даёт улучшение ECE с 0,070 до 0,043 (в производственную конфигурацию не включена).
    \item \textbf{Слой 3 — байесовский валидатор.} Сеть (структура — Hill Climb + BIC) вычисляет $P_B(\hat{y}_k \mid \mathbf{e})$; ниже порога $\theta_k$ — флаг и передача дальше. Роль пересмотрена (3.1): не классификатор, а валидатор плотностей.
    \item \textbf{Слой 4 — запасной слой на большой языковой модели.} Взаимозаменяемые модели семейств GPT \cite{ref_41} и Claude \cite{ref_42}: Sonnet 4.5, GPT-4o, Gemini 2.5 Flash, открытая gpt-oss-120b, малая llama-3.2-3b (сопоставление — 3.3.2). Структурированный запрос со схемой и партнёр-доступными полями; возвращает значение или честный отказ.
\end{itemize}

Слои упорядочены по росту удельной стоимости решения: регулярные выражения почти бесплатны, ML работает на уже вычисленном векторном представлении, байесовский вывод требует одного прохода по сети, запасной слой — самый дорогой.

Структура обученной байесовской сети для категории \texttt{chocolate} приведена на рисунке~\ref{fig:bayes_dag_ch2}. Узлы графа — целевые атрибуты и бренд; направленные рёбра показывают условные зависимости, найденные алгоритмом Hill Climb на серебряной выборке. Бренд — корневой узел; через производный признак \texttt{brand\_has\_organic\_marker} он влияет на \texttt{chocolate\_type} и \texttt{is\_organic}. Тип шоколада задаёт процент какао и подтип \texttt{chocolate\_extra}, от которого расходятся рёбра к \texttt{contains\_nuts}, \texttt{nutri\_score\_grade}, \texttt{protein\_class} и повторно к \texttt{is\_organic} — это согласуется с интуицией о связи органичности, бренда и подтипа. Эмпирическая роль сети как точечного валидатора и количественные характеристики приведены в 3.3.4.

\begin{figure}
  \includegraphics[width=0.8\textwidth]{fig_3_5_bayes_dag.png}
  \caption{Структура байесовской сети категории \texttt{chocolate} (направленный ациклический граф; алгоритм Hill Climb с критерием BIC, библиотека pgmpy 1.1.2)}
  \label{fig:bayes_dag_ch2}
\end{figure}

\subsubsection{Обоснование выбора компонентов стека}

\textbf{Каскад вместо сквозного решения.} Сквозной all-LLM даёт 83--94~\% точности (3.3.2), но стоимость линейно растёт с числом ячеек; каскад при равной или большей точности экономит порядка двух порядков за счёт ступенчатого распределения нагрузки. Сквозной all-ML не подходит — у части атрибутов значения вычисляются из данных за пределами товарного описания, к тому же есть сценарий холодного старта новых категорий.

\textbf{SBERT \texttt{paraphrase-multilingual-mpnet-base-v2}} \cite{ref_20} выбран за бесплатную многоязычность (50+ языков в едином пространстве без дообучения), умеренную 768-мерную размерность и готовность к работе без доменной адаптации. Архитектура MPNet объединяет преимущества маскированного и перестановочного языкового моделирования и устойчиво обгоняет дистиллированные MiniLM-варианты на задачах семантической близости.

\textbf{XGBoost} \cite{ref_14} с регуляризацией (subsample=0.8, colsample=0.8, ранняя остановка, gamma=0.1) выбран как Слой 2 за устойчивость на плотных векторных признаках при ограниченной разметке (n$\approx$1250 на категорию), калибруемость вероятностей (изотоническая регрессия \cite{ref_28}) и миллисекундный вывод на CPU. Поэтапная оптимизация по остаткам и явные L1/L2/$\gamma$-штрафы дают более тонкие нелинейные границы и более прямой контроль переобучения, чем у Random Forest~\cite{ref_53,ref_54}; на независимых табличных стендах \cite{ref_55,ref_56} градиентный бустинг устойчиво обходит Random Forest и нейросетевые архитектуры в характерной для Слоя 2 размерности. Эмпирическое сравнение (\S\,\ref{sec:method-comparison}) подтверждает выбор: при сопоставимой точности XGBoost даёт наименьшую ECE (0,026 против 0,031 у MLP и 0,035 у LogReg).

\textbf{Байесовская сеть на pgmpy \cite{ref_35} 1.1.2.} По сравнению с MRF и автоэнкодерами дискретная байесовская сеть допускает построение структуры графа (Hill Climb + BIC \cite{ref_36}) с интерпретируемым ориентированным графом, дешёвый вывод через Variable Elimination и явное выражение условных зависимостей; pgmpy поддерживает и построение структуры, и оценку параметров через \texttt{BayesianEstimator}.

\textbf{Шлюз API на Go при ML-сервисе на Python.} Демо-система — два сервиса: шлюз на Go (\texttt{:8080}) и ML-сервис на Python/FastAPI (\texttt{:8001}). Go подходит для сетевого ввода-вывода (малое потребление памяти, быстрый запуск, конкурентный I/O); Python-экосистема (PyTorch, sentence-transformers, XGBoost, pgmpy) безальтернативна для ML-инференса. Разделение соответствует типовой производственной архитектуре и позволяет масштабировать компоненты независимо.

\textbf{Открытая LLM как запасной слой.} Конфигурация многомодельна через OpenRouter (Claude Sonnet 4.5, GPT-4o, Gemini 2.5 Flash, открытая gpt-oss-120b, малая llama-3.2-3b). Открытую gpt-oss-120b можно запустить в инфраструктуре заказчика без передачи данных вовне; коммерческие LLM взаимозаменяемы, что снижает зависимость от провайдера; каскад компенсирует слабость малой модели — при самостоятельном применении gpt-oss-120b проигрывает Sonnet 14--15~п.~п., но в составе каскада превосходит её безусловное использование на 21,3~п.~п. (3.3.2).

\subsubsection{Стек технологий — сводная таблица}

Сводное представление принятых решений приведено в таблице 2.1.

\begin{xltabular}{\textwidth}{|>{\raggedright\arraybackslash}p{3.0cm}|>{\raggedright\arraybackslash}p{5.5cm}|>{\raggedright\arraybackslash}X|}
\caption{Сводная таблица стека технологий с обоснованием выбора} \\ \hline
    Компонент & Технология & Ключевое обоснование \\ \hline
\endfirsthead
\multicolumn{3}{l}{Продолжение таблицы} \\ \hline
    Компонент & Технология & Ключевое обоснование \\ \hline
\endhead
    Шлюз API & Go, стандартная библиотека net/http & Малое потребление ресурсов, быстрый запуск, отсутствие внешних зависимостей \\ \hline
    ML-сервис & Python 3.14, FastAPI, Uvicorn & Доступ к ML-экосистеме, автоматическая валидация через Pydantic, поддержка OpenAPI \\ \hline
    Векторное представление & SBERT \texttt{paraphrase-\allowbreak multilingual-\allowbreak mpnet-\allowbreak base-v2} & Многоязычность (50+ языков), 768-мерный вектор (архитектура MPNet), отсутствие необходимости в дообучении \\ \hline
    Классификатор Слой 2 & XGBoost с регуляризацией; перекрёстная изотоническая калибровка на эталонной выборке снижает среднюю ECE с 0,070 до 0,043 (3.3.3.3), внедрение в производственную конфигурацию вынесено в раздел «Перспективы развития темы» заключения & Устойчивость на плотных признаках, калибруемые вероятности, дешёвый вывод \\ \hline
    Байесовская сеть & pgmpy 1.1.2, \texttt{Discrete}\allowbreak\texttt{BayesianNetwork} & Структурное обучение Hill Climb + BIC, дешёвый вывод (Variable Elimination), интерпретируемый граф \\ \hline
    Категорайзер Слой 0 & LightGBM (поверх TF-IDF-представления) + порог OOD & Многоклассовая классификация с вероятностями + явный отказ при низкой уверенности \\ \hline
    Запасной слой & OpenRouter API: Sonnet 4.5, GPT-4o, Gemini 2.5 Flash, gpt-oss-120b, llama-3.2-3b & Единый интерфейс, возможность локального развёртывания, независимость от поставщика \\ \hline
    Фронтенд & HTML/CSS/JavaScript без сборочной системы & Минимальные требования к окружению, прозрачность демонстрации \\ \hline
\end{xltabular}

\subsection{Сценарий работы программного обеспечения}

\subsubsection{Сквозной пример прохождения запроса}

Рассмотрим прохождение на примере «Spaghetti \#5 Barilla» (повторяется в 3.1.2). Партнёр шлёт \texttt{POST /api/enrich} с полями category=null, product\_name=«Spaghetti \#5 Barilla», brands=«Barilla», ingredients\_text=«Durum wheat semolina, water», quantity=«500g». Поле category пустое — категоризация на системе.

\textbf{Шаг 1. Определение категории (Слой 0).} Возвращает $c = \text{pasta}$ с уверенностью $0{,}97$; OOD не сработал, выбрана схема $\mathcal{A}_{\text{pasta}}$ (8 атрибутов).

\textbf{Шаг 2. Поочерёдная обработка атрибутов.} Каскад обходит атрибуты по схеме. Для pasta\_shape срабатывает шаблон Слоя 1 на «spaghetti» с уверенностью $1{,}0$ — следующие слои не вызываются. Для grain\_type Слой 1 не даёт совпадения; Слой 2 выдаёт \texttt{"wheat"} с $P_{\text{ML}} = 0{,}99 \geq \tau = 0{,}5$, валидатор подтверждает ($P_B = 0{,}85 \geq \theta = 0{,}04$) — отметка ml. Для nutri\_score\_grade — класс \texttt{"A"}, $P_{\text{ML}} = 0{,}83$, валидатор подтверждает. Для is\_filled $P_{\text{ML}} = 0{,}45 < \tau$ — ячейка идёт в Слой 4.

\textbf{Шаг 3. Формирование ответа.} ML-сервис собирает поатрибутный ответ с агрегированными счётчиками и возвращает JSON шлюзу (формат — 3.1.2). В этом сценарии 7 из 8 атрибутов закрыты дешёвыми слоями (1 regex + 6 ML), 1 — Слоем 4. Именно такое распределение и даёт каскаду выигрыш в матрице «точность — стоимость» (3.3.2).

\subsection{Выводы по главе 2}

\begin{enumerate}[label=\arabic*)]
    \item Задача обогащения формализована как поатрибутное предсказание $\hat{y}_k$ с явным отказом ($\hat{y}_k = \text{null}$); каскадная логика — последовательная фильтрация по уверенности с конъюнктивным условием $P_{\text{ML}} \geq \tau_k \wedge P_B \geq \theta_k$ для принятия предсказания Слоя 2 как ансамблевая оценка уверенности.
    \item Архитектурный выбор каждой компоненты (каскад вместо сквозного решения; SBERT, XGBoost, pgmpy; Go-шлюз + Python-ML; открытая gpt-oss-120b) обоснован эксплуатационно, методологически и стратегически. Это превращает слабую открытую модель в работоспособную часть производственного стека.
    \item Сквозной сценарий подтверждает: большинство атрибутов закрывается дешёвыми слоями, обращение к LLM точечно — отсюда и доминирование каскада по Парето над одиночным вызовом большой языковой модели (подтверждается в 3.3).
\end{enumerate}
